% 4.5 章 · 多重网格与快速求解
% 这一章只回答一个问题：为什么「几次扫描 + 几层网格」能把误差压下去，
% 以及这句话在什么条件下是错的。四步链条：
%   1) 平滑子是逐频率的：高频一次掉到 1/3 以下，低频几乎不动（放大率 1-O(h^2)）；
%   2) 粗层校正把低频交给粗网格：一维插值 + 满权限制下 Galerkin 恒等式 R A_h P = A_H 精确成立；
%   3) 两者复合的两网格算子谱恰为 {0, 1/9}（= 平滑因子的平方），
%      V/W 循环的代价分别是 O(n_0) 与 O(n_0 log n_0)；
%   4) 非线性必须走 FAS：把它当线性问题解，会收敛到「线性方程的解」这个假解。
% 本章所有结论都与 numcore/pde/multigrid.py 的实现一一对应，
% 对应关系集中写在「本章代码与结论的对应」一节。
% 断言附表由 tools/build_book.py 从实际运行结果生成，\input 进来，不手抄。

\numlabchapter{多重网格与快速求解}{4.5}

\section{这一章要解决什么}

1.5 章把迭代法的收敛速度归结到一件事上：迭代矩阵的谱。对一维 Poisson 模型问题，
加权 Jacobi 与 Gauss--Seidel 的谱半径都是
$\rho(T)=1-O(h^2)$（定理~\ref{thm:mg-spectrum}、\ref{lem:mg-stationary}），
于是把误差压下一个固定倍数要 $\Theta(h^{-2})$ 次扫描：$h$ 每减半，扫描数翻两番。
CG 好一些，但 $\kappa(A_h)=O(h^{-2})$，Chebyshev 界给出的迭代数仍是 $O(h^{-1})$
（定理~\ref{thm:mg-mgcg}）。**这两条路都不可能做到 $O(1)$ 次迭代。**

多重网格的出发点是：误差不是「一个数」，而是**一组频率**，
而不同频率该交给不同的网格去处理。

\begin{enumerate}
  \item \textbf{高频（振荡）分量}：几次加权 Jacobi 就能把它按固定倍数压下去，
        因为平滑子的放大率在频率半带 $\theta\ge\pi/2$ 上一致地小
        （$\omega=2/3$ 时不超过 $1/3$，定理~\ref{thm:mg-smooth}）。
        这是\textbf{平滑}，不是「迭代收敛」——低频在这个判据里几乎不动。
  \item \textbf{低频（光滑）分量}：它在细网格上「看起来像」粗网格上的一条曲线，
        所以在粗网格上它是高频，可以用同样的平滑子加粗层校正处理。
        把这一步递归下去就是 V/W 循环（\S\ref{sec:mg-cycle}）。
  \item \textbf{两者的复合}：$\nu_1$ 次前扫描 + 粗层校正 + $\nu_2$ 次后扫描。
        在一维模型问题上，$\nu_1=\nu_2=1$、$\omega=2/3$ 的两网格算子谱恰好是
        $\{0,\,1/9\}$（定理~\ref{thm:mg-tg-spectrum}）——
        因子与 $h$ 无关，这就是「网格无关」的全部内容。
  \item \textbf{代价}：每层视野内的点数按 $2^{-l}$ 缩小，所以
        V 循环一个周期的代价是 $O(n_0)$、W 循环是 $O(n_0\log n_0)$
        （定理~\ref{thm:mg-cost}），加上 $O(1)$ 的周期数，总代价 $O(n_0)$。
  \item \textbf{边界条件}：线性校正算子必须是一个**固定**的线性算子，
        这样才能当 CG 的预条件子用（\S\ref{sec:mg-mgcg}）；
        非线性问题必须用 FAS，否则会收敛到不是解的假解（\S\ref{sec:mg-fas}）。
\end{enumerate}

\subsection*{记号（与 \texttt{numcore/pde/multigrid.py} 完全一致）}

模型问题是**一维 Dirichlet Poisson 方程**
\begin{equation}
  -u''(x)=f(x),\quad x\in(0,1),\qquad u(0)=u(1)=0,
  \label{eq:mg-model}
\end{equation}
在网格 $x_i=ih$、$h=1/(n+1)$ 上离散为
\begin{equation}
  A_h=h^{-2}\operatorname{tridiag}(-1,2,-1)\in\mathbb R^{n\times n},
  \qquad A_hu_h=f_h .
  \label{eq:mg-A}
\end{equation}
$A_h$ 与 1.5 章的 $\operatorname{laplacian\_1d}(n)$ 只差一个因子 $h^{-2}$。

\textbf{网格层级}：$n_l$ 是第 $l$ 层的内点数，$l=0$ 最细，
\begin{equation}
  n_l=\frac{n_0+1}{2^{\,l}}-1,\qquad
  h_l=\frac{1}{n_l+1}=2^{\,l}h_0,\qquad H=2h ,
  \label{eq:mg-hierarchy}
\end{equation}
最粗层只有 $1$ 个内点。第 $l$ 层的算子就是 \eqref{eq:mg-A} 在粗网格上**重新离散**：
$A_H=H^{-2}\operatorname{tridiag}(-1,2,-1)$。粗点 $j$ 与细点重合：在一基下标下，
粗点 $j$ 落在细点 $2j$ 上（代码里数组下标从 $0$ 起，那里的写法是粗点 $j$ 落在细点 $2j+1$ 上）。

\textbf{转移算子}：线性插值 $P$（粗 $\to$ 细）与满权限制 $R$（细 $\to$ 粗）
\begin{equation}
  (Pu)_{2j}=u_j,\qquad (Pu)_{2j\pm1}=\tfrac12\bigl(u_j+u_{j\pm1}\bigr),
  \qquad
  (Rr)_j=\tfrac14\bigl(r_{2j-1}+2r_{2j}+r_{2j+1}\bigr),
  \label{eq:mg-PR}
\end{equation}
越界的粗值取 $0$（即 Dirichlet 边界）。

\textbf{平滑子}：$A=D-L-U$（$D$ 对角、$L,U\ge0$ 严格三角），
\begin{equation}
  T_\omega=I-\omega D^{-1}A\quad(\text{加权 Jacobi}),\qquad
  T_{\mathrm{GS}}=(D-L)^{-1}U\quad(\text{前向 Gauss--Seidel}),
  \label{eq:mg-smoothers}
\end{equation}
$T_{\mathrm{GS}}$ 就是 $\omega=1$ 的 SOR：$T_{\mathrm{GS}}=(D-L)^{-1}\bigl((1-\omega)D+\omega U\bigr)\big|_{\omega=1}$。

\textbf{模态}：$\theta_k=k\pi/(n+1)=k\pi h$，$k=1,\dots,n$；
正弦向量 $v_k(i)=\sin(i\theta_k)$。频率半带指 $\theta_k\ge\pi/2$，即 $k\ge(n+1)/2$。
逐模态放大率
\begin{equation}
  \lambda_k(A_h)=4h^{-2}\sin^2\frac{\theta_k}{2},\qquad
  \lambda_k(D^{-1}A_h)=1-\cos\theta_k,\qquad
  \mu_k(\omega)=1-\omega\,(1-\cos\theta_k).
  \label{eq:mg-modes}
\end{equation}
记 $s_k=\sin^2(\theta_k/2)$、$c_k=\cos^2(\theta_k/2)=1-s_k$；$N=n+1$，于是 $v_{N-k}(i)=(-1)^{i+1}v_k(i)$。

\textbf{算子}：$E_{\rm TG}$ 是两网格算子的误差传播算子，$E_l$ 是第 $l$ 层循环算子，
$B$ 是循环作为预条件子时「残差 $\mapsto$ 修正」的算子（\S\ref{sec:mg-mgcg}）。

\section{模型问题、平滑子与转移算子}
\label{sec:mg-setting}

\begin{theorem}[$A_h$ 的谱与加权 Jacobi 的逐模态放大率]
\label{thm:mg-spectrum}
$\{v_k\}_{k=1}^{n}$ 是 $\mathbb R^n$ 的正交基，并且
\begin{equation}
  A_hv_k=\lambda_k(A_h)v_k=4h^{-2}\sin^2\frac{\theta_k}{2}\,v_k,
  \qquad
  D^{-1}A_hv_k=(1-\cos\theta_k)v_k
  \label{eq:mg-sine}
\end{equation}
（$D=2h^{-2}I$）。因此加权 Jacobi 的迭代矩阵 $T_\omega=I-\omega D^{-1}A_h$
在同一组基下对角：
\begin{equation}
  T_\omega v_k=\mu_k(\omega)v_k,\qquad
  \mu_k(\omega)=1-\omega(1-\cos\theta_k)=1-2\omega\sin^2\frac{\theta_k}{2},
  \label{eq:mg-mu}
\end{equation}
它的第 $k$ 个模态误差系数**精确**按 $\mu_k$ 衰减（没有模态混合）：
$c_k^{(m)}=\mu_k^m c_k^{(0)}$。此外 $T_\omega$ 对称，且关于 $A_h$-内积自伴：
$A_hT_\omega=T_\omega^{\mathsf T}A_h$。
\end{theorem}

\begin{proof}
第一式是 1.5 章 一维 Laplacian 特征分解的直接改写，这里给出自洽的验算。
由积化和差
\begin{equation}
  \sin((i-1)\theta)+\sin((i+1)\theta)=2\cos\theta\,\sin(i\theta)
\end{equation}
得 $\bigl(\operatorname{tridiag}(-1,2,-1)v_k\bigr)_i=2(1-\cos\theta_k)\sin(i\theta_k)
=4\sin^2(\theta_k/2)\sin(i\theta_k)$，乘上 $h^{-2}$ 即得第一式。
正交性由 $\sum_{i=1}^{n}\sin(i\theta_j)\sin(i\theta_k)=\frac{N}{2}\delta_{jk}$
（$N=n+1$）给出。$D=2h^{-2}I$ 是纯量矩阵，于是
$D^{-1}A_hv_k=\frac12\cdot 4\sin^2(\theta_k/2)v_k=(1-\cos\theta_k)v_k$，
第二式得证。

$T_\omega=I-\omega D^{-1}A_h$ 是 $D^{-1}A_h$ 的多项式（一次），
所以与它同特征向量，特征值 \eqref{eq:mg-mu}。
在正弦基下 $T_\omega$ 是对角矩阵 $\operatorname{diag}(\mu_1,\dots,\mu_n)$，
而该基正交，故 $c_k^{(m)}=\mu_k^mc_k^{(0)}$ 精确成立
（这一点与 Gauss--Seidel 形成对照：后者在同一组基下**不是**对角的，见 \S\ref{sec:mg-gs}）。

对称性：$D^{-1}A_h$ 对称（$A_h$ 对称、$D$ 是纯量矩阵），故 $T_\omega$ 对称。
关于 $A_h$-内积自伴：$A_hT_\omega=A_h-\omega A_hD^{-1}A_h$，
而 $A_hD^{-1}A_h$ 对称，故 $A_hT_\omega$ 对称，即
$A_hT_\omega=(A_hT_\omega)^{\mathsf T}=T_\omega^{\mathsf T}A_h$。
\end{proof}

\begin{theorem}[最优平滑权重与平滑因子]
\label{thm:mg-smooth}
对任意 $\omega>0$，平滑子在高频半带上的最坏放大率是
\begin{equation}
  \mu_{\rm smooth}(\omega):=\max_{\theta_k\ge\pi/2}\bigl|\mu_k(\omega)\bigr|
  =\max\bigl(|1-\omega|,\;|1-2\omega|\bigr),
  \label{eq:mg-smoothmax}
\end{equation}
两个端点分别对应 $\theta=\pi/2$（$s_k=1/2$）与 $\theta=\pi$（$s_k=1$）。
它在 $\omega=2/3$ 取到最小值 $1/3$，即
\begin{equation}
  \omega_\ast=\frac23,\qquad
  \mu_{\rm smooth}\Bigl(\frac23\Bigr)=\frac13,
  \qquad
  \Bigl|\mu_k\Bigl(\frac23\Bigr)\Bigr|=\Bigl|1-\frac43\sin^2\frac{\theta_k}{2}\Bigr|\le\frac13
  \quad(\theta_k\ge\pi/2).
  \label{eq:mg-smoothing-factor}
\end{equation}
\end{theorem}

\begin{proof}
把 $\mu_k=1-2\omega s$ 看成 $s=\sin^2(\theta_k/2)$ 的一次函数。
$\theta_k\ge\pi/2$ 时 $s\in[1/2,1]$，一次函数在区间上的最大模在端点取到，于是
\begin{equation}
  \mu_{\rm smooth}(\omega)=\max\Bigl(\bigl|1-2\omega\cdot\tfrac12\bigr|,\;
  \bigl|1-2\omega\cdot1\bigr|\Bigr)=\max(|1-\omega|,|1-2\omega|).
\end{equation}
下面分四段求这个最大值的最小值（$g(\omega)=\max(|1-\omega|,|1-2\omega|)$）：
\begin{itemize}
  \item $\omega\le\frac12$：$g=1-\omega\ge\frac12>\frac13$；
  \item $\frac12\le\omega\le\frac23$：$2\omega-1\le1-\omega$，故 $g=1-\omega\ge\frac13$；
  \item $\frac23\le\omega\le1$：$g=2\omega-1\ge\frac13$；
  \item $\omega\ge1$：$g\ge|1-2\omega|=2\omega-1\ge1>\frac13$。
\end{itemize}
所以最小值是 $1/3$，且只在 $\omega=2/3$ 取到；代回即得
$\mu_k(2/3)=1-\frac43\sin^2(\theta_k/2)$，
在 $s=1/2$ 与 $s=1$ 处分别等于 $1/3$ 与 $-1/3$，模都是 $1/3$。
\end{proof}

\begin{corollary}[「低频几乎不动」的定量含义]
\label{cor:mg-lowmode}
最低模态 $\theta_1=\pi h$ 的放大率满足
\begin{equation}
  \mu_1\Bigl(\frac23\Bigr)=1-\frac23\bigl(1-\cos(\pi h)\bigr)
  =1-\frac{\pi^2h^2}{3}+O(h^4)
  =1-\Theta(h^2),
  \label{eq:mg-lowmode}
\end{equation}
而同一权重下任何高频模态都不超过 $1/3$。因此：
\begin{enumerate}
  \item 「低频几乎不动」的精确含义是\textbf{放大率 $=1-\Theta(h^2)$}：
        $h$ 减半，距离 $1$ 的距离缩小到四分之一；
        单靠平滑子把一个低频分量压下一个固定倍数需要 $\Theta(h^{-2})$ 次扫描
        （首项系数由 $1-\mu_1\approx\pi^2h^2/3$ 给出）。
  \item 高频分量每扫描一次至少缩小 $3$ 倍：$\Theta(1)$ 次扫描就够。
        这就是「平滑」与「收敛」的分工：平滑子不是求解器，它只负责高频。
\end{enumerate}
\end{corollary}

\begin{proof}
由 $\cos x=1-\frac{x^2}{2}+\frac{x^4}{24}+O(x^6)$，
\begin{equation}
  \frac23\bigl(1-\cos(\pi h)\bigr)
  =\frac23\Bigl(\frac{\pi^2h^2}{2}-\frac{\pi^4h^4}{24}+O(h^6)\Bigr)
  =\frac{\pi^2h^2}{3}+O(h^4),
\end{equation}
故 \eqref{eq:mg-lowmode} 成立；$1-\mu_1=\frac{\pi^2h^2}{3}(1+O(h^2))$ 给出下界与上界同阶，
即 $\Theta(h^2)$，于是把误差压下固定倍数所需的扫描数是
$\Theta\bigl(1/(1-\mu_1)\bigr)=\Theta(h^{-2})$。
高频部分由定理~\ref{thm:mg-smooth} 的 \eqref{eq:mg-smoothing-factor} 直接给出：
$\theta_k\ge\pi/2$ 时 $|\mu_k|\le1/3$，$m$ 次扫描后不超过 $3^{-m}$。
\end{proof}

\begin{definition}[两个网格算子]
\label{def:mg-tg}
粗层校正算子与两网格误差传播算子定义为
\begin{equation}
  C:=I-PA_H^{-1}RA_h,\qquad
  E_{\rm TG}:=S^{\nu_2}\,C\,S^{\nu_1},\qquad S:=T_{2/3},
  \label{eq:mg-tg}
\end{equation}
其中 $\nu_1$ 次前扫描、$\nu_2$ 次后扫描，粗层用精确解（$A_H^{-1}$）。
\end{definition}

\begin{theorem}[$R=\frac12P^{\mathsf T}$ 与一维 Galerkin 恒等式]
\label{thm:mg-galerkin}
按 \eqref{eq:mg-PR} 定义的转移算子满足
\begin{equation}
  R=\tfrac12P^{\mathsf T},
  \qquad
  P^{\mathsf T}A_hP=2A_H,
  \qquad
  \boxed{\,RA_hP=A_H\,}.
  \label{eq:mg-galerkin}
\end{equation}
即：一维的 Galerkin 粗层算子与在粗网格上重新离散得到的 $A_H$ **逐元素相同**，
缺陷恒为 $0$（不是「小量」）。
\end{theorem}

\begin{proof}
\emph{第一步：$R=\frac12P^{\mathsf T}$。}
$P$ 的第 $j$ 列在一基下标下有：第 $2j$ 个分量 $1$，第 $2j\pm1$ 个分量 $1/2$。
故
\begin{equation}
  (P^{\mathsf T}r)_j=r_{2j}+\tfrac12\bigl(r_{2j-1}+r_{2j+1}\bigr),
  \qquad
  \tfrac12(P^{\mathsf T}r)_j=\tfrac14\bigl(r_{2j-1}+2r_{2j}+r_{2j+1}\bigr)=(Rr)_j .
\end{equation}

\emph{第二步：$P^{\mathsf T}A_hP=2A_H$。}
设 $u$ 是粗层向量，$y=Pu$。线性插值的两个关键取值：
\begin{equation}
  y_{2j}=u_j,\qquad y_{2j\pm1}=\tfrac12\bigl(u_j+u_{j\pm1}\bigr).
\end{equation}
先算 $A_hy$ 的两类分量。在粗点上，
\begin{equation}
  (A_hy)_{2j}=h^{-2}\bigl(2u_j-\tfrac12(u_{j-1}+u_j)-\tfrac12(u_j+u_{j+1})\bigr)
  =h^{-2}\bigl(u_j-\tfrac12u_{j-1}-\tfrac12u_{j+1}\bigr);
\end{equation}
在中点上，因为插值是\textbf{线性}的，二阶差分正好为零：
\begin{equation}
  (A_hy)_{2j+1}=h^{-2}\bigl(2\cdot\tfrac12(u_j+u_{j+1})-u_j-u_{j+1}\bigr)=0 .
\end{equation}
于是
\begin{equation}
  (P^{\mathsf T}A_hPu)_j=(A_hy)_{2j}+\tfrac12\bigl((A_hy)_{2j-1}+(A_hy)_{2j+1}\bigr)
  =h^{-2}\bigl(u_j-\tfrac12u_{j-1}-\tfrac12u_{j+1}\bigr).
\end{equation}
用 $h^{-2}=4H^{-2}$ 整理：
\begin{equation}
  h^{-2}\bigl(u_j-\tfrac12u_{j-1}-\tfrac12u_{j+1}\bigr)
  =2H^{-2}\bigl(2u_j-u_{j-1}-u_{j+1}\bigr)=2(A_Hu)_j .
\end{equation}
第三步把第一步代入第二步即得 $RA_hP=\frac12P^{\mathsf T}A_hP=A_H$。
\end{proof}

\begin{corollary}[缩放常数被恒等式定死]
\label{cor:mg-scaling}
若把限制算子换成 $R=cP^{\mathsf T}$（$c\ne\frac12$），则
$RA_hP=2cA_H$：粗层算子的缩放与限制算子的缩放是一个**配对**。
只有 $c=\frac12$ 时 $RA_hP=A_H$ 才成立；$c=1$（即 $R=P^{\mathsf T}$）
等价于把粗层校正放大一倍，$c$ 取得更小则是欠校正。
\end{corollary}

\begin{proof}
$RA_hP=cP^{\mathsf T}A_hP=2cA_H$（定理~\ref{thm:mg-galerkin} 第二步）。
$2c=1$ 当且仅当 $c=1/2$。
\end{proof}

\begin{remark}[二维：Galerkin 与重新离散不是同一个算子]
一维里 $P^{\mathsf T}A_hP=2A_H$ 是**恒等式**，所以本章代码可以在每一层直接用
$\operatorname{poisson\_1d}(n_l)$ 当粗层算子，而不必真的去算 $P^{\mathsf T}A_hP$。
二维没有这个运气：双线性插值配满权限制（$R=\frac14Q^{\mathsf T}$）时，
Galerkin 粗层算子是\textbf{九点格式}（四个角点权重非零），
而不是五点重新离散 $\operatorname{tridiag}_2$。
两者相差一个 $O(h^2)$ 的扰动项，于是二维必须**二选一**：
\begin{itemize}
  \item Galerkin（变分）路线：$A_H=RA_hP$，$R=\frac12P^{\mathsf T}$，
        粗层校正是 $A_h$-正交投影，\S\ref{sec:mg-mgcg} 的对称性证明逐字成立，但粗层是九点、更贵；
  \item 重新离散路线：粗层是五点、更便宜，但 $RA_hP\ne A_H$，
        粗层校正不再是投影，「精确解在 $\operatorname{range}(P)$ 上被一次消掉」的结论也不再有。
\end{itemize}
一维的特殊之处正在于两条路线重合。
\end{remark}

\section{两网格算子：投影、h 无关性与精确谱}
\label{sec:mg-tg}

误差传播的解释：设真解为 $u^\ast$，误差 $e=u-u^\ast$。一次两网格循环里
前扫描把误差变成 $S^{\nu_1}e$，粗层校正再把它变成 $C S^{\nu_1}e$，
后扫描给出 $E_{\rm TG}e=S^{\nu_2}CS^{\nu_1}e$。逐项看：
\begin{itemize}
  \item $S$ 负责高频（定理~\ref{thm:mg-smooth}）；
  \item $C=I-PA_H^{-1}RA_h$ 负责低频：它把误差限制到粗网格、
        在粗网格上解出修正、再插值回来。
\end{itemize}

\begin{lemma}[粗层校正是投影，且在 $\operatorname{range}(P)$ 上精确]
\label{lem:mg-cgc-projection}
设 $M:=PA_H^{-1}RA_h$。则
\begin{enumerate}
  \item $M^2=M$，$C=I-M$ 是投影；$\operatorname{range}(M)=\operatorname{range}(P)$，
        $\ker M=\ker(RA_h)=\ker(P^{\mathsf T}A_h)$；
  \item 对任意粗层向量 $w$ 有 $C(Pw)=0$：粗层校正把落在
        $\operatorname{range}(P)$ 里的误差**一次消干净**（这是「修正方程」意义上的精确性）；
  \item 反过来，$Cx=x$ 对一切 $x\in\ker(RA_h)$ 成立：粗网格「看不见」的误差分量
        粗层校正一点也帮不上忙，只能交给平滑子；
  \item $C$ 关于 $A_h$-内积自伴：$A_hC$ 对称（即 $C$ 是 $A_h$-正交投影的补投影）。
\end{enumerate}
\end{lemma}

\begin{proof}
(1) 用 $RA_hP=A_H$（定理~\ref{thm:mg-galerkin}）：
\begin{equation}
  M^2=PA_H^{-1}\bigl(RA_hP\bigr)A_H^{-1}RA_h=PA_H^{-1}A_HA_H^{-1}RA_h=M .
\end{equation}
故 $C^2=I-2M+M^2=I-M=C$。又 $\operatorname{range}(M)\subseteq\operatorname{range}(P)$，
且 $MPw=PA_H^{-1}RA_hPw=Pw$，故 $\operatorname{range}(M)=\operatorname{range}(P)$；
$\ker M=\ker(RA_h)=\ker(P^{\mathsf T}A_h)$（$R=\frac12P^{\mathsf T}$，$A_h$ 可逆）。
(2) $C(Pw)=Pw-PA_H^{-1}RA_hPw=Pw-PA_H^{-1}A_Hw=0$。
(3) $x\in\ker(RA_h)$ 时 $Mx=0$，故 $Cx=x$。
(4) $A_hM=A_hPA_H^{-1}RA_h=A_hP\bigl(\tfrac12P^{\mathsf T}A_hP\bigr)^{-1}\tfrac12P^{\mathsf T}A_h
=A_hP(P^{\mathsf T}A_hP)^{-1}P^{\mathsf T}A_h$，这是对称矩阵（转置不变），
于是 $A_hC=A_h-A_hM$ 也对称。
\end{proof}

\begin{theorem}[平滑性 $+$ 逼近性 $\Rightarrow$ $h$ 无关]
\label{thm:mg-two-property}
设存在与 $h$ 无关的常数 $\eta(\nu_1),C$，使得对一切 $h$：
\begin{enumerate}
  \item \textbf{平滑性}（per-mode 版本的算子形式）：
        $\|A_hS^{\nu_1}\|\le\eta(\nu_1)h^{-2}$；
  \item \textbf{逼近性}（变分版本：粗层空间确实能逼近细层解）：
        $\|A_h^{-1}-PA_H^{-1}R\|\le Ch^{2}$；
  \item 后扫描不放大：$\|S^{\nu_2}\|\le1$。
\end{enumerate}
则两网格算子的范数与 $h$ 无关：
\begin{equation}
  \|E_{\rm TG}\|\le\eta(\nu_1)\,C .
  \label{eq:mg-two-property}
\end{equation}
\end{theorem}

\begin{proof}
关键是把粗层校正写成「逼近误差 $\times\,A_h$」的形式：
\begin{equation}
  I-PA_H^{-1}RA_h=\bigl(A_h^{-1}-PA_H^{-1}R\bigr)A_h .
\end{equation}
代入 \eqref{eq:mg-tg}：
\begin{equation}
  E_{\rm TG}=S^{\nu_2}\bigl(A_h^{-1}-PA_H^{-1}R\bigr)\,A_hS^{\nu_1},
\end{equation}
取范数并用次可乘性：
\begin{equation}
  \|E_{\rm TG}\|\le\|S^{\nu_2}\|\cdot\|A_h^{-1}-PA_H^{-1}R\|\cdot\|A_hS^{\nu_1}\|
  \le 1\cdot Ch^{2}\cdot\eta(\nu_1)h^{-2}=\eta(\nu_1)C .
\end{equation}
$h^{-2}$ 被 $h^{2}$ 精确抵消——这就是「两条性质合起来给网格无关性」的全部机制。
\end{proof}

\begin{remark}
两条性质缺一不可：只有平滑性（高频掉得快）而没有逼近性，低频永远留在原地；
只有逼近性（粗层能表示低频）而没有平滑性，粗网格看不见的高频分量无人处理。
上面这条界是\textbf{安全但不锐}的：它只保证 $h$ 无关，常数乘积一般远大于真值。
本章模型问题的锐结论是下一个定理：谱半径恰为 $1/9$，与 $h$ 无关。
\end{remark}

\begin{theorem}[两网格算子的精确谱：$\sigma(E_{\rm TG})=\{0,\frac19\}$]
\label{thm:mg-tg-spectrum}
考虑一维模型问题，$\omega=2/3$ 的加权 Jacobi 平滑子 $S=I-\frac13\operatorname{tridiag}(-1,2,-1)$，
$\nu_1=\nu_2=1$，$R=\frac12P^{\mathsf T}$，$A_H=RA_hP$（一维时即重新离散，定理~\ref{thm:mg-galerkin}）。
记 $N=n+1$，$\theta_k=k\pi/N$，$s_k=\sin^2(\theta_k/2)$，$c_k=\cos^2(\theta_k/2)$。则
\begin{enumerate}
  \item 粗层校正对单个模态的作用是**混合一对模态**：
        \begin{equation}
          Cv_k=s_k\bigl(v_k+v_{N-k}\bigr),\qquad
          Cv_{N-k}=c_k\bigl(v_k+v_{N-k}\bigr)
          \label{eq:mg-C-pair}
        \end{equation}
        （$k=N/2$ 时两式重合，给出 $Cv_{N/2}=v_{N/2}$）；
  \item 每个二维不变子空间 $X_k=\operatorname{span}\{v_k,v_{N-k}\}$ 在 $E_{\rm TG}$ 下不变，
        且 $E_{\rm TG}|_{X_k}$ 的特征值是 $0$ 与
        \begin{equation}
          \lambda_k=s_k\mu_k^2+c_k\mu_{N-k}^2,
          \qquad \mu_j=1-\tfrac23(1-\cos\theta_j);
          \label{eq:mg-lambda-k}
        \end{equation}
  \item 恒等式 $\lambda_k=\dfrac19$ 对每个 $k$ 成立。因此
        \begin{equation}
          \sigma(E_{\rm TG})=\Bigl\{0,\ \frac19\Bigr\},
          \qquad
          \rho(E_{\rm TG})=\frac19,
          \qquad
          E_{\rm TG}^2=\frac19E_{\rm TG}\ne0 .
          \label{eq:mg-rho-tg}
        \end{equation}
        $1/9$ 的重数是 $\lceil N/2\rceil$（$N$ 偶数时为 $N/2$，即每对模态贡献一个），
        特征向量是 $S(v_k+v_{N-k})$。
\end{enumerate}
\end{theorem}

\begin{proof}
\emph{(1) 的第一步：$w_k:=v_k+v_{N-k}\in\ker(P^{\mathsf T}A_h)$。}
记 $s_i=\sin(i\theta_k)$。由 $v_{N-k}(i)=(-1)^{i+1}v_k(i)$ 得
\begin{equation}
  w_k(i)=\begin{cases}2s_i,& i\ \text{奇},\\ 0,& i\ \text{偶}.\end{cases}
\end{equation}
于是
\begin{equation}
  (A_hw_k)_i=\begin{cases}
    2h^{-2}w_k(i)=4h^{-2}s_i, & i\ \text{奇},\\[2pt]
    -h^{-2}\bigl(w_k(i-1)+w_k(i+1)\bigr)=-4h^{-2}\cos\theta_k\,s_i, & i\ \text{偶}.
  \end{cases}
\end{equation}
（偶数点用了 $s_{i-1}+s_{i+1}=2\cos\theta_k s_i$。）粗层点是一基下标里的偶数点，故
\begin{equation}
  \begin{aligned}
    (P^{\mathsf T}A_hw_k)_j
    &=(A_hw_k)_{2j}+\tfrac12\bigl((A_hw_k)_{2j-1}+(A_hw_k)_{2j+1}\bigr)\\
    &=-4h^{-2}\cos\theta_k\,s_{2j}+\tfrac12\cdot 4h^{-2}\bigl(s_{2j-1}+s_{2j+1}\bigr)=0 .
  \end{aligned}
\end{equation}
（最后一步同样用 $s_{2j-1}+s_{2j+1}=2\cos\theta_k s_{2j}$。）

\emph{(1) 的第二步：$v_k-s_kw_k\in\operatorname{range}(P)$。}
粗层点上 $w_k$ 为零，中点上是 $2s_i$，所以
\begin{equation}
  (v_k-s_kw_k)_{2j}=s_{2j},\qquad
  (v_k-s_kw_k)_{2j\pm1}=(1-2s_k)s_{2j\pm1}=\cos\theta_k\,s_{2j\pm1}.
\end{equation}
而 $\cos\theta_k s_{2j\pm1}=\frac12\bigl(s_{2j}+s_{2j\pm2}\bigr)=\frac12\bigl(s_{2j}+s_{2(j\pm1)}\bigr)$，
正是线性插值公式；边界上 $s_{2(n_c+1)}=\sin(N\theta_k)=\sin(k\pi)=0$，
与 Dirichlet 的越界取零一致。故 $v_k-s_kw_k=Pu$，
其中粗向量 $u_j=s_{2j}=\sin(j\cdot2\theta_k)$。

\emph{(1) 的结论。} 由引理~\ref{lem:mg-cgc-projection}，$C$ 是「沿 $\operatorname{range}(P)$
投影到 $\ker(RA_h)$ 上」的投影：$C$ 在 $\ker(RA_h)$ 上是恒等、在 $\operatorname{range}(P)$ 上是零。
把 $v_k=s_kw_k+(v_k-s_kw_k)$ 分解成
$\ker(RA_h)\ni s_kw_k$ 与 $\operatorname{range}(P)\ni v_k-s_kw_k$ 两部分，
作用 $C$ 后只剩第一项：$Cv_k=s_kw_k$。
同法（把 $k$ 换成 $N-k$，并用 $s_{N-k}=c_k$）得 $Cv_{N-k}=c_kw_k$。
$k=N/2$ 时 $v_{N-k}=v_k$、$s_k=1/2$，两式都退化为 $Cv_k=v_k$。

\emph{(2)} $S$ 在正弦基下对角（定理~\ref{thm:mg-spectrum}），
$C$ 把 $X_k$ 映到 $\operatorname{span}\{w_k\}\subset X_k$，
所以 $X_k$ 在 $E_{\rm TG}=SCS$ 下不变。在基 $(v_k,v_{N-k})$ 下
\begin{equation}
  S=\begin{pmatrix}\mu_k&0\\0&\mu_{N-k}\end{pmatrix},\qquad
  C=\begin{pmatrix}s_k&c_k\\ s_k&c_k\end{pmatrix},
\end{equation}
\begin{equation}
  E_{\rm TG}=SCS=\begin{pmatrix}
    s_k\mu_k^2 & c_k\mu_k\mu_{N-k}\\[2pt]
    s_k\mu_k\mu_{N-k} & c_k\mu_{N-k}^2
  \end{pmatrix}.
\end{equation}
迹是 $\lambda_k$，行列式是 $s_kc_k\mu_k^2\mu_{N-k}^2-c_ks_k\mu_k^2\mu_{N-k}^2=0$，
故特征值为 $0$ 与 $\lambda_k$。

\emph{(3)} 记 $c=\cos\theta_k$。由 $\mu_j=1-\frac23(1-\cos\theta_j)=\frac13(1+2\cos\theta_j)$、
$\cos\theta_{N-k}=-\cos\theta_k=-c$ 得
$\mu_k=\frac13(1+2c)$、$\mu_{N-k}=\frac13(1-2c)$，又 $s_k=\frac{1-c}{2}$、$c_k=\frac{1+c}{2}$，于是
\begin{equation}
  \lambda_k=\frac{(1-c)(1+2c)^2+(1+c)(1-2c)^2}{18} .
\end{equation}
展开分子：$(1-c)(1+4c+4c^2)=1+3c-4c^3$，
$(1+c)(1-4c+4c^2)=1-3c+4c^3$，两者相加恰为 $2$，
故 $\lambda_k=\frac{2}{18}=\frac19$。

\emph{(4) 全部特征值与 $E^2=\frac19E$。}
$X_k$（$k=1,\dots,\lceil N/2\rceil$）两两正交、直和是整个 $\mathbb R^n$，
所以 $\sigma(E_{\rm TG})=\{0,\frac19\}$，$1/9$ 每个 $k$ 一个、共 $\lceil N/2\rceil$ 个。
在每个不变子空间上特征多项式是 $\lambda^2-\frac19\lambda$（一维的 $k=N/2$ 情形是 $\lambda-\frac19$），
由 Cayley--Hamilton，$E_{\rm TG}$ 在每个 $X_k$ 上满足 $E^2=\frac19E$，故整体成立，特别地 $\rho=1/9$。
特征向量：把 $S$ 作用在 $w_k$ 上，$Sw_k=\mu_kv_k+\mu_{N-k}v_{N-k}$，于是
\begin{equation}
  E_{\rm TG}\bigl(Sw_k\bigr)=SCS^2w_k=SC\bigl(\mu_k^2v_k+\mu_{N-k}^2v_{N-k}\bigr)
  =S\bigl(s_k\mu_k^2+c_k\mu_{N-k}^2\bigr)w_k=\tfrac19\,Sw_k .
\end{equation}
\end{proof}

\begin{remark}[为什么是「平方」，以及 $1/9$ 不是「每步都缩到 $1/9$」]
\begin{enumerate}
  \item 平滑子出现\textbf{两次}（前扫描一次、后扫描一次），故 \eqref{eq:mg-lambda-k} 里是 $\mu^2$；
        权重 $s_k,c_k$ 来自插值/限制的 $\frac14[1,2,1]$ 模板。
        对一般的 $\omega$，$\lambda_k(\omega)=s_k\mu_k(\omega)^2+c_k\mu_{N-k}(\omega)^2$
        依赖于 $k$；只有 $\omega=2/3$ 时恒等式
        $\frac{1-c}{2}(1+2c)^2+\frac{1+c}{2}(1-2c)^2=2$ 把它压成常数，
        于是 $\rho_{\rm TG}=\mu_{\rm smooth}^2=(1/3)^2=1/9$。
  \item $E_{\rm TG}^2=\frac19E_{\rm TG}$ 说明：第一个循环之后误差就落进
        $\operatorname{range}(E_{\rm TG})$，此后每个循环精确乘 $1/9$——不会「更快」。
  \item $1/9$ 是\textbf{谱半径}，即渐近速率。$E_{\rm TG}$ 并不欧氏对称
        （$C$ 只关于 $A_h$-内积自伴），所以单步下降由 $\|E_{\rm TG}\|$ 控制，
        一般比 $1/9$ 慢；把 $1/9$ 读成「每步误差除以 9」是过度解读。
\end{enumerate}
\end{remark}

\section{V 循环与 W 循环：递归与代价}
\label{sec:mg-cycle}

粗层不可能真的用 $A_H^{-1}$：那要一次精确求解。实际做法是递归——
在粗层再跑一个（更小的）多重网格循环，直到最粗层（只有 $1$ 个内点）用固定次数的扫描。
\textbf{固定次数}这一点是本质的：它保证整个循环是一个固定的（仿射）线性算子，
既可以用行列式讨论非奇异性（\S\ref{sec:mg-cycle}），也可以当 CG 的预条件子（\S\ref{sec:mg-mgcg}）。

\begin{definition}[V/W 循环的递归]
\label{def:mg-cycle}
设层数 $L$，$E_l$ 是第 $l$ 层循环的误差传播算子，$\gamma\in\{1,2\}$（$\gamma=1$ 是 V 循环、
$\gamma=2$ 是 W 循环），最粗层做 $m$ 次固定扫描：
\begin{equation}
  E_{L-1}=S_{L-1}^{m},
  \qquad
  E_l=S_l^{\nu_2}\Bigl[I-P_l\bigl(I-E_{l+1}^{\gamma}\bigr)A_{l+1}^{-1}R_lA_l\Bigr]S_l^{\nu_1},
  \quad l<L-1 .
  \label{eq:mg-recursion}
\end{equation}
\end{definition}

递归 \eqref{eq:mg-recursion} 的来源：线性循环从零初值出发时是仿射的，
$x\mapsto E_{l+1}x+(I-E_{l+1})A_{l+1}^{-1}r$，从零出发跑 $\gamma$ 次得到粗层响应
\begin{equation}
  \sigma_\gamma(r)=\Bigl(I+E_{l+1}+\cdots+E_{l+1}^{\gamma-1}\Bigr)\bigl(I-E_{l+1}\bigr)A_{l+1}^{-1}r
  =\bigl(I-E_{l+1}^{\gamma}\bigr)A_{l+1}^{-1}r .
  \label{eq:mg-coarse-response}
\end{equation}
把它代进「粗层校正 $x\mapsto x+P\sigma(Rr)$」即得 \eqref{eq:mg-recursion}。
$E_{l+1}=0$（粗层精确求解）时它退化为两网格形式 \eqref{eq:mg-tg}。

\begin{theorem}[V 循环 $O(n_0)$、W 循环 $O(n_0\log n_0)$]
\label{thm:mg-cost}
以一个「最细层等价点扫描」为代价单位（在第 $l$ 层做一次扫描的代价是 $n_l/n_0$）。
设 $n_l=(n_0+1)/2^l-1\le n_02^{-l}$，每层 $\nu_1+\nu_2$ 次扫描，最粗层 $m$ 次（$m$ 是常数）。
记 $W_V(l)$、$W_W(l)$ 为第 $l$ 层一个循环的代价（以点扫描为单位），则
\begin{equation}
  W_V(l)=(\nu_1+\nu_2)n_l+W_V(l+1),\qquad W_V(L-1)=m\,n_{L-1},
\end{equation}
\begin{equation}
  W_W(l)=(\nu_1+\nu_2)n_l+2W_W(l+1),\qquad W_W(L-1)=m\,n_{L-1}.
\end{equation}
解出来是
\begin{equation}
  W_V(0)=(\nu_1+\nu_2)\sum_{l=0}^{L-2}n_l+m\,n_{L-1}
  <\Bigl(2(\nu_1+\nu_2)+m\,2^{-(L-1)}\Bigr)n_0=O(n_0),
  \label{eq:mg-cost-v}
\end{equation}
\begin{equation}
  W_W(0)=(\nu_1+\nu_2)\sum_{l=0}^{L-2}2^{\,l}\,n_l+2^{\,L-1}m\,n_{L-1}
  =(\nu_1+\nu_2)(L-1)n_0+m\,n_0=O\bigl(n_0\log n_0\bigr).
  \label{eq:mg-cost-w}
\end{equation}
\end{theorem}

\begin{proof}
\emph{V 循环（几何级数）。}
逐层展开递推：$W_V(0)=(\nu_1+\nu_2)\sum_{l=0}^{L-2}n_l+m\,n_{L-1}$。
用 $n_l\le n_02^{-l}$ 与几何级数 $\sum_{l\ge0}2^{-l}=2$：
\begin{equation}
  \sum_{l=0}^{L-2}n_l\le n_0\sum_{l=0}^{L-2}2^{-l}<2n_0,
  \qquad
  m\,n_{L-1}\le m\,n_02^{-(L-1)} .
\end{equation}
每层只跑一个粗层循环，所以层数只以\textbf{加性}方式进入：整个 V 循环的代价被最细层支配，
$W_V(0)=O(n_0)$，与层数 $L$ 无关（$L$ 只影响常数 $2(\nu_1+\nu_2)+m2^{-(L-1)}$）。

\emph{W 循环（项数 $\times$ 每项代价）。}
递推 $W_W(l)=(\nu_1+\nu_2)n_l+2W_W(l+1)$ 展开一次就多一倍：
\begin{equation}
  W_W(0)=(\nu_1+\nu_2)\sum_{l=0}^{L-2}2^{\,l}n_l+2^{\,L-1}m\,n_{L-1}.
\end{equation}
第 $l$ 层被访问 $2^{\,l}$ 次，每次代价 $n_l$，乘积 $2^{\,l}n_l\le n_02^{\,l}2^{-l}=n_0$
\textbf{与 $l$ 无关}——这正是 W 循环比 V 循环贵一个对数的原因：
\begin{equation}
  W_W(0)\le(\nu_1+\nu_2)(L-1)n_0+m\,n_0 .
\end{equation}
嵌套网格给出 $n_0+1=2^{\,L-1}(n_{L-1}+1)$，故 $L-1=\log_2\frac{n_0+1}{2}\le\log_2n_0$，
于是 $W_W(0)=O(n_0\log n_0)$。
\end{proof}

\begin{remark}
转移算子（限制与插值）每层代价也是 $O(n_l)$，把它们算进去结论不变：
V 循环 $O(n_0)$、W 循环 $O(n_0\log n_0)$。本章代码的计数器
（\texttt{mg\_apply} 的 \texttt{counter}）只累计点扫描数，即 \eqref{eq:mg-cost-v}、\eqref{eq:mg-cost-w}
里的 $W_V,W_W$；转移代价按上面的道理同阶，不改变结论。
\end{remark}

\begin{theorem}[V 循环的因子不是零]
\label{thm:mg-nonzero}
沿用 \eqref{eq:mg-recursion} 的记号。
\begin{enumerate}
  \item \textbf{（粗层精确时）} $E_l$ 退化为两网格算子，由定理~\ref{thm:mg-tg-spectrum}
        它的谱是 $\{0,1/9\}$：$\rho=1/9>0$，且 $E_{\rm TG}^2=\frac19E_{\rm TG}\ne0$，
        故 $E_{\rm TG}$ 不是幂零的；它在高频子空间上也不为零——
        例如 $w_k=v_k+v_{N-k}$ 满足 $E_{\rm TG}(Sw_k)=\frac19Sw_k\ne0$。
  \item \textbf{（一般递归）} 与两网格算子的差是「粗层误差被重新注入」这一项：
        \begin{equation}
          E_l-E_l^{\rm TG}=S_l^{\nu_2}P_l\,E_{l+1}^{\gamma}A_{l+1}^{-1}R_lA_l\,S_l^{\nu_1},
          \label{eq:mg-cycle-diff}
        \end{equation}
        它为零当且仅当 $E_{l+1}=0$，即当且仅当粗层被精确求解。
        粗层只要是用固定次数扫描做的，$E_{l+1}\ne0$，循环算子与两网格算子就\textbf{不是}同一个算子。
  \item \textbf{（非零下界）} 由 Sylvester 行列式恒等式
        $\det(I-XY)=\det(I-YX)$ 与 Galerkin 恒等式 $R_lA_lP_l=A_{l+1}$：
        \begin{equation}
          \begin{aligned}
            \det E_l&=\det(S_l)^{\nu_1+\nu_2}\cdot
            \det\bigl(I-P_l(I-E_{l+1}^{\gamma})A_{l+1}^{-1}R_lA_l\bigr)\\
            &=\det(S_l)^{\nu_1+\nu_2}\cdot\det\bigl(E_{l+1}\bigr)^{\gamma}.
          \end{aligned}
          \label{eq:mg-det}
        \end{equation}
        于是逐层递推得
        $\det E_0=\Bigl[\prod_{l=0}^{L-2}\det(S_l)^{\nu_1+\nu_2}\Bigr]\det(S_{L-1})^{m}$。
        只要每层的平滑子与最粗层算子都非奇异，$E_0$ 就非奇异，
        从而\textbf{不是幂零的}：$\rho(E_0)>0$，V 循环的因子是一个正数。
\end{enumerate}
\end{theorem}

\begin{proof}
(1) 见定理~\ref{thm:mg-tg-spectrum}（取 $\gamma=1$、$E_{l+1}=0$ 时
\eqref{eq:mg-recursion} 与 \eqref{eq:mg-tg} 一致）。
(2) 把 $I-P(I-E^{\gamma})A^{-1}RA=I-PA^{-1}RA+PE^{\gamma}A^{-1}RA$ 代入 \eqref{eq:mg-recursion}，
与 $E_l^{\rm TG}=S^{\nu_2}(I-PA_{l+1}^{-1}R_lA_l)S^{\nu_1}$ 相减即得 \eqref{eq:mg-cycle-diff}。
(3) 记 $X=P_l(I-E_{l+1}^{\gamma})A_{l+1}^{-1}$、$Y=R_lA_l$，则
\begin{equation}
  \begin{aligned}
    \det\bigl(I-XY\bigr)&=\det\bigl(I_{n_{l+1}}-YX\bigr)
    =\det\Bigl(I-R_lA_lP_l\bigl(I-E_{l+1}^{\gamma}\bigr)A_{l+1}^{-1}\Bigr)\\
    &=\det\bigl(I-\bigl(I-E_{l+1}^{\gamma}\bigr)\bigr)=\det\bigl(E_{l+1}^{\gamma}\bigr),
  \end{aligned}
\end{equation}
第二个等号用 Sylvester，第三个用 $R_lA_lP_l=A_{l+1}$。
（$X$ 是 $n_l\times n_{l+1}$、$Y$ 是 $n_{l+1}\times n_l$，两者行列式因子同。）
幂零矩阵的特征值全为零、行列式为零，故行列式非零即排除了幂零性，$\rho>0$。
\end{proof}

\begin{remark}[应该期待什么]
$\omega=2/3$、$\nu_1=\nu_2=1$ 的 V 循环在一维模型问题上的因子是 $O(0.1)$ 量级——
与两网格的 $1/9$ 同阶、略大，因为粗层没有精确求解（\eqref{eq:mg-cycle-diff} 的注入项是正的）。
把误差压到 $10^{-8}$ 需要的循环数由 $\rho^k\le10^{-8}$ 定出：
$k\ge \ln(10^{8})/\ln(1/\rho)$；取理想值 $\rho=1/9$ 得 $k\approx8.4$，即 $9$ 个循环
（这是 $1/9$ 的直接推算，不是实测值；实际因子略大，循环数也略多）。
W 循环的因子更接近 $1/9$，代价是每周期 $O(n_0\log n_0)$（定理~\ref{thm:mg-cost}）。
\end{remark}

\section{多重网格作为预条件子（MGCG）}
\label{sec:mg-mgcg}

把循环当求解器用，收敛速率由 $\rho(E)$ 决定；把它当\textbf{预条件子}用，
则是让 CG 去处理谱分布。PCG 算法（与 \texttt{numcore/linalg/krylov.py} 逐行一致）：
\begin{equation}
  \begin{aligned}
    &r_0=b-Ax_0,\quad z_0=M^{-1}r_0,\quad p_0=z_0,\\
    &\alpha_k=\frac{r_k^{\mathsf T}z_k}{p_k^{\mathsf T}Ap_k},\quad
    x_{k+1}=x_k+\alpha_kp_k,\quad
    r_{k+1}=r_k-\alpha_kAp_k,\\
    &z_{k+1}=M^{-1}r_{k+1},\quad
    \beta_k=\frac{r_{k+1}^{\mathsf T}z_{k+1}}{r_k^{\mathsf T}z_k},\quad
    p_{k+1}=z_{k+1}+\beta_kp_k .
  \end{aligned}
  \label{eq:mg-pcg}
\end{equation}
在多重网格里 $M^{-1}r$ 就是「以 $r$ 为右端、零初值跑一个循环」，
所以需要回答：这个算子是不是一个合格的预条件子？

\begin{theorem}[对称循环算子是对称正定的]
\label{thm:mg-cycle-spd}
设 $\nu_1=\nu_2=\nu$，平滑子 $S=S^{\mathsf T}$ 对称且与 $D^{-1}$ 交换，
$R=\frac12P^{\mathsf T}$，粗层算子是 Galerkin 的（一维即重新离散），
粗层响应 $\sigma$ 对称半正定（三种典型情形都满足：精确解 $\sigma=A_H^{-1}$；
粗层 $m$ 次扫描 $\sigma=\omega\sum_{j=0}^{m-1}S_{\!c}^{\,j}D_c^{-1}$；
递归的粗层循环响应——对层数归纳）。
则循环算子（零初值、右端 $r$ 的一次循环，记 $z=Br$）有闭式
\begin{equation}
  \boxed{\;B=\omega\sum_{j=0}^{2\nu-1}S^{\,j}D^{-1}\;+\;S^{\nu}P\,\sigma\,R\,S^{\nu}\;}
  \label{eq:mg-B-closed}
\end{equation}
两项都欧氏对称，第一项正定、第二项半正定，故 $B=B^{\mathsf T}>0$。
\end{theorem}

\begin{proof}
记一次循环的三步（零初值、右端 $r$）：前扫描
\begin{equation}
  x_1=\omega\sum_{j=0}^{\nu-1}S^{\,j}D^{-1}r=:\Sigma_\nu r,
\end{equation}
（因为单次扫描是 $x\mapsto Sx+\omega D^{-1}r$，$\nu$ 次求和即得；
也等价于 $x_1=(I-S^{\nu})A^{-1}r$），粗层校正
\begin{equation}
  x_2=x_1+P\sigma R\bigl(r-Ax_1\bigr),
\end{equation}
后扫描
\begin{equation}
  x_3=S^{\nu}x_2+\Sigma_\nu r .
\end{equation}
代入并用 $\Sigma_\nu=\omega\sum_{j<\nu}S^jD^{-1}$、$A\Sigma_\nu=I-S^{\nu}$
（由 $I-S=\omega D^{-1}A$ 与 $A$ 和 $D^{-1}$、$S$ 都交换得到）：
\begin{equation}
  B=S^{\nu}\bigl(I-P\sigma RA\bigr)\Sigma_\nu+S^{\nu}P\sigma R+\Sigma_\nu
  =\Sigma_{2\nu}+S^{\nu}P\sigma RS^{\nu},
\end{equation}
最后一步用了 $I-A\Sigma_\nu=S^{\nu}$。

\emph{对称性}：$S$ 与 $D^{-1}$ 都对称且交换，故 $S^{\,j}D^{-1}$ 对称（写成
$S^{\lfloor j/2\rfloor}D^{-1}S^{\lceil j/2\rceil}$），第一项对称。
第二项里 $(P\sigma R)^{\mathsf T}=R^{\mathsf T}\sigma^{\mathsf T}P^{\mathsf T}
=\frac12P\sigma P^{\mathsf T}=P\sigma\frac12P^{\mathsf T}=P\sigma R$ 对称，再被对称的 $S^{\nu}$ 夹住仍对称。
\emph{正定性}：第一项是正定矩阵 $D^{-1}$ 与半正定矩阵的乘积之和（$\omega>0$、$S^{\,j}D^{-1}$ 半正定），
第二项半正定，故 $B$ 正定。
\end{proof}

\begin{theorem}[MGCG 保持 CG 的最优性；迭代数 $h$ 无关]
\label{thm:mg-mgcg}
设 $B$ 如 \eqref{eq:mg-B-closed} 对称正定，取 $M=B^{-1}$ 代入 \eqref{eq:mg-pcg}。则
\begin{enumerate}
  \item PCG 良定义，且保持 CG 的最优性：第 $k$ 步迭代在
        $x_0+\mathcal K_k(BA,Br_0)$ 中最小化 A-范数误差，并有 Chebyshev 界
        \begin{equation}
          \|e_k\|_A\le2\left(\frac{\sqrt\kappa-1}{\sqrt\kappa+1}\right)^{k}\|e_0\|_A,
          \qquad \kappa=\kappa(BA)=\frac{\lambda_{\max}(BA)}{\lambda_{\min}(BA)};
          \label{eq:mg-cheb}
        \end{equation}
  \item 因为 $B$ 是一个固定的循环算子，$\lambda(BA)$ 落在一个\textbf{与 $h$ 无关}的区间里，
        所以迭代数 $O(1)$；而普通 CG 的 $\kappa(A_h)=O(h^{-2})$，
        \eqref{eq:mg-cheb} 给出 $k=O(h^{-1})$。这正是 MGCG 与 CG 的分水岭。
\end{enumerate}
\end{theorem}

\begin{proof}
$B$ 对称正定，故 $M=B^{-1}$ 对称正定，\eqref{eq:mg-pcg} 就是标准的预条件共轭梯度法，
其最优性与 \eqref{eq:mg-cheb} 的推导与 1.5 章 定理（CG 最优性与 Chebyshev 界）逐字相同，
只需把那里的 $M$ 换成这里的 $B^{-1}$：$\kappa=\kappa(M^{-1}A)=\kappa(BA)$。

第二条分两步。
\emph{普通 CG}：由 \eqref{eq:mg-modes}，
\begin{equation}
  \kappa(A_h)=\frac{\lambda_n(A_h)}{\lambda_1(A_h)}
  =\frac{\sin^2(\theta_n/2)}{\sin^2(\theta_1/2)}
  =\frac{\cos^2(\pi h/2)}{\sin^2(\pi h/2)}=\cot^2\frac{\pi h}{2}=O(h^{-2}),
\end{equation}
（用了 $\theta_1=\pi h$、$\theta_n=\pi-\pi h$）。于是 $\sqrt\kappa=\cot\frac{\pi h}{2}\approx\frac{2}{\pi h}$，
\eqref{eq:mg-cheb} 里的收缩因子
$\frac{\sqrt\kappa-1}{\sqrt\kappa+1}\approx1-\frac{2}{\sqrt\kappa}\approx1-\pi h$，
把误差压下一个固定倍数需要 $O(1/(\pi h))=O(h^{-1})$ 步。
\emph{MGCG}：$B$ 由 \eqref{eq:mg-B-closed} 给出，其中 $S,D^{-1},P,\sigma,R$ 全是固定的算子组合，
$\sigma$ 的最粗层部分只是固定次数的扫描，故 $BA$ 与 $h$ 无关（至多差一个与 $h$ 无关的相似变换），
$\kappa(BA)=O(1)$，\eqref{eq:mg-cheb} 给出 $k=O(1)$。
\end{proof}

\begin{remark}[「作为求解器」与「作为预条件子」的代价之别]
\begin{itemize}
  \item \textbf{作为求解器}：每次迭代 $=$ 一个循环，代价 $O(n_0)$（V）或 $O(n_0\log n_0)$（W）；
        迭代数 $O(1)$（由 $\rho(E)$ 定）；总代价 $O(n_0)$（V）或 $O(n_0\log n_0)$（W）。
        实现最短、每步最省，但收敛速率完全由 $\rho(E)$ 决定，
        对 $\omega$、$\nu$、循环类型、限制方式都敏感。
  \item \textbf{作为预条件子}：每次 PCG 迭代 $=$ 一个循环 $+$ 若干次 $O(n_0)$ 的向量运算
        （内积、$axpy$、方向更新），单步常数更大（本章的实现里一次 V 循环约合几个最细层等价代价）；
        迭代数由 $\kappa(BA)$ 定，也是 $O(1)$，通常比循环数更少，
        而且 CG 的最优性让方法对参数不敏感：谱的聚簇被 Krylov 空间自动利用。
  \item \textbf{共同前提}：预条件子必须是\textbf{固定}的线性算子。
        最粗层的扫描次数必须固定（本章 \texttt{COARSE\_SWEEPS}），
        不能按残差停机——按残差停机时循环是数据依赖的非线性算子，
        $B$ 不再固定，\eqref{eq:mg-cheb} 的最优性证明失效。这也是
        \S\ref{sec:mg-cycle} 一开始强调「固定次数扫描」的第二个理由。
\end{itemize}
\end{remark}

\section{非线性问题：FAS}
\label{sec:mg-fas}

模型问题换成
\begin{equation}
  -u''(x)+u(x)^3=f(x),\qquad x\in(0,1),\qquad u(0)=u(1)=0,
  \label{eq:mg-nl}
\end{equation}
离散为 $A_h(u)=f_h$，其中
\begin{equation}
  \bigl(A_h(u)\bigr)_i=h^{-2}\bigl(2u_i-u_{i-1}-u_{i+1}\bigr)+u_i^3 .
  \label{eq:mg-nl-op}
\end{equation}
平滑子不再是线性扫描，而是逐点的标量非线性解：
把 $u_{i-1},u_{i+1}$ 固定，对 $u_i$ 解
\begin{equation}
  h^2u_i^3+2u_i=u_{i-1}+u_{i+1}+h^2f_i .
  \label{eq:mg-cubic}
\end{equation}

\begin{theorem}[标量三次方程恰有一个实根；二分区间无条件有效]
\label{thm:mg-cubic}
对任意 $s\in\mathbb R$，方程
\begin{equation}
  g(u):=h^2u^3+2u=s
  \label{eq:mg-cubic-root}
\end{equation}
在 $\mathbb R$ 上恰有一个实根；而且这个根落在区间
$\bigl[\min(0,s),\max(0,s)\bigr]$ 内。因此在该区间上做二分法无条件收敛
（不需要任何 Newton 型收敛假设）。此外，$\omega=1$ 时平滑子的不动点与细层非线性方程
$A_h(u)=f_h$ **完全等价**。
\end{theorem}

\begin{proof}
$g'(u)=3h^2u^2+2>0$，故 $g$ 严格单调增；又 $g(u)\to\pm\infty$（$u\to\pm\infty$），
所以对每个 $s$ 恰有一个实根。
设 $s\ge0$：若 $u<0$ 则 $g(u)<g(0)=0\le s$，无根；
若 $u>s\ge0$ 则 $g(u)=h^2u^3+2u>2u>2s\ge s$，也无根。
所以根必在 $[0,s]=[\min(0,s),\max(0,s)]$ 内（$s<0$ 时对称地得 $[s,0]$）。
二分法在这段区间上每一步都把区间减半，$52$ 步后区间长度小于 $2^{-52}$，已是机器精度。

不动点等价性：一次逐点解把 $u_i$ 换成 \eqref{eq:mg-cubic} 的根，
故 $u$ 是不动点当且仅当 \eqref{eq:mg-cubic} 对每个 $i$ 成立；
把它除以 $h^2$ 并移项得
$h^{-2}(2u_i-u_{i-1}-u_{i+1})+u_i^3=f_i$，即 $A_h(u)=f_h$。
\end{proof}

\begin{theorem}[FAS：粗层方程是「被限制的细层方程」，且没有假不动点]
\label{thm:mg-fas}
记 $R=\frac12P^{\mathsf T}$、$\tilde u$ 为前扫描后的迭代。FAS 的粗层问题是
\begin{equation}
  A_H(u_H)=f_H(\tilde u):=R\bigl(f_h-A_h(\tilde u)\bigr)+A_H\bigl(R\tilde u\bigr),
  \label{eq:mg-fas-coarse}
\end{equation}
初值 $u_H^{(0)}=R\tilde u$，修正为 $u\leftarrow\tilde u+P(u_H-R\tilde u)$。
则
\begin{enumerate}
  \item \textbf{（一致性/推导）} 若 $A_h(u)=f_h$，则 \eqref{eq:mg-fas-coarse} 的右端恰为
        $A_H(Ru)$，即 $u_H=Ru$ 精确满足粗层方程；粗层问题因此是细层问题的**限制**，
        不含额外的截断项。（这正是必须保留 $A_H(R\tilde u)$ 这一项的原因。）
  \item \textbf{（无假不动点）} 取 $\nu_1=\nu_2=1$、平滑子为 \eqref{eq:mg-cubic} 的精确逐点解、
        粗层精确求解。若 $u$ 是 FAS 循环的不动点，则 $A_h(u)=f_h$。
\end{enumerate}
\end{theorem}

\begin{proof}
\textbf{(1)} 把 $A_h(u)=f_h$ 代入：$f_H=R(f_h-f_h)+A_H(Ru)=A_H(Ru)$，
于是 $u_H=Ru$ 是粗层方程的解；反过来，粗层方程两边同减 $A_H(R\tilde u)$ 后右端只剩限制残差，
所以在解 $u$ 处没有任何 $O(1)$ 的偏差项。

\textbf{(2)} 记 $w=\tilde u+d$、$d=P(u_H-R\tilde u)$、$e=u-\tilde u$。
\begin{itemize}
  \item \emph{平滑恒等式}：若 $v=S(z)$ 是逐点解，则
        $h^{-2}(2v_i-z_{i-1}-z_{i+1})+v_i^3=f_i$，与 $(A_h(v))_i$ 的定义相减得
        \begin{equation}
          A_h\bigl(S(z)\bigr)=f_h+h^{-2}(L+U)\bigl(z-S(z)\bigr),
          \label{eq:mg-smooth-identity}
        \end{equation}
        其中 $\bigl((L+U)x\bigr)_i=x_{i-1}+x_{i+1}$。特别地 $S(z)=z$ 时 $A_h(z)=f_h$。
  \item \emph{后扫描}（$z=w$、$v=u$）：$\rho(u):=f_h-A_h(u)=-h^{-2}(L+U)(w-u)$。
  \item \emph{前扫描}（$z=u$、$v=\tilde u$）：$\rho(\tilde u)=f_h-A_h(\tilde u)=-h^{-2}(L+U)e$。
  \item \emph{中值定理}：$A_h(u)-A_h(\tilde u)=De$，其中
        $D=\int_0^1A_h'(\tilde u+te)\,dt$ 对称正定，因为
        $A_h'(v)=h^{-2}\operatorname{tridiag}(-1,2,-1)+3\operatorname{diag}(v^2)$。
        与前两条相减得 $De=h^{-2}(L+U)(d-e)$，即
        \begin{equation}
          \begin{aligned}
            &\Lambda e=h^{-2}(L+U)d,\\
            &\Lambda:=D+h^{-2}(L+U)=3\operatorname{diag}\!\Bigl(\int_0^1v^2\Bigr)+2h^{-2}I
            \quad\text{（对角、正定）}.
          \end{aligned}
          \label{eq:mg-lambda}
        \end{equation}
        （$D$ 的次对角 $-h^{-2}$ 与 $(L+U)$ 的 $+h^{-2}$ 正好抵消，$\Lambda$ 是对角的。）
  \item \emph{粗层方程}：由精确求解，$A_H(u_H)-A_H(R\tilde u)=R\rho(\tilde u)=-Rh^{-2}(L+U)e$。
        再用一次中值定理（粗层）：$A_H(u_H)-A_H(R\tilde u)=M_Hc$，
        $c=u_H-R\tilde u$，$M_H$ 对称正定。代入 \eqref{eq:mg-lambda} 与 $d=Pc$：
        \begin{equation}
          \Bigl[M_H+\tfrac12P^{\mathsf T}h^{-2}(L+U)\Lambda^{-1}h^{-2}(L+U)P\Bigr]c=0 .
        \end{equation}
        方括号是正定矩阵（$M_H$ 正定）加半正定矩阵（$P^{\mathsf T}(\cdot)P$），故正定、可逆，
        于是 $c=0$；回到 $d=Pc=0$ 与 \eqref{eq:mg-lambda} 得 $e=0$，
        即 $u=\tilde u$ 是平滑子的不动点，由定理~\ref{thm:mg-cubic} 得 $A_h(u)=f_h$。
\end{itemize}
\end{proof}

\begin{remark}[粗层扫描次数有限时]
实现里最粗层只做固定次数 $m$ 的扫描（\texttt{COARSE\_SWEEPS}），粗层方程留有缺陷
$\delta_H=f_H-A_H(u_H)$。上面的第五步变成
$\bigl[M_H+\frac12P^{\mathsf T}(\cdot)P\bigr]c=-\delta_H$，
于是不动点的细层残差被 $\|\delta_H\|$ 线性控制：粗层解得越准，假不动点越接近真解；
只有粗层精确（或逐层归纳到各层自身的不动点）时「无假不动点」才是严格等式。
\end{remark}

\begin{theorem}[把非线性当线性解会收敛到错方程的解]
\label{thm:mg-false}
若在残差与粗层算子中都丢掉三次项（即用线性算子
$A_h^{\rm lin}=h^{-2}\operatorname{tridiag}(-1,2,-1)$ 代替 $A_h(\cdot)$，
平滑子也退化成线性扫描），则整个循环就是**线性**多重网格的校正格式，
其不动点恰是线性方程组 $A_h^{\rm lin}u=f_h$ 的解；
在这个点上真非线性残差等于 $u^3$，不会随迭代变小——收敛到的是解错方程得到的假解。
\end{theorem}

\begin{proof}
丢掉三次项后循环是仿射的：$u\mapsto Gu+c$，$G$ 就是线性多重网格的误差传播算子。
由定理~\ref{thm:mg-tg-spectrum}（或其递归版本，只要因子 $<1$），$\rho(G)<1$，
故 $I-G$ 可逆，不动点唯一，且由线性理论它就是
$u^\ast=(A_h^{\rm lin})^{-1}f_h$（把 $u^\ast$ 代进仿射格式，误差项 $Gu^\ast+c-u^\ast$ 为零）。
在 $u^\ast$ 处 $A_h(u^\ast)-f_h=(u^\ast)^3$，一般不为零；
两个残差（线性残差 $\approx0$ 与真残差 $\approx u^3$）的比值可以任意大。
\end{proof}

\begin{remark}[丢掉 $A_H(R\tilde u)$ 也一样致命]
若把 \eqref{eq:mg-fas-coarse} 换成 $A_H(u_H)=R(f_h-A_h(\tilde u))$（即丢掉 $A_H(R\tilde u)$），
并把 $u_H$ 当修正用，则不动点（修正为零 $\Rightarrow u_H=R\tilde u$）只满足
\begin{equation}
  R\bigl(f_h-A_h(u)\bigr)=A_H(Ru),
\end{equation}
这是**限制后的残差方程**，不是细层方程：$R$ 只有 $n_c$ 行，
$R(f_h-A_h(u))=0$ 有大量非零解，格式因此有假不动点。
FAS 的写法把 $A_H(R\tilde u)$ 留在粗层方程里，正是为了让不动点方程回到
$R(f_h-A_h(u))=0$ 之前的「方程限制」形式，再由定理~\ref{thm:mg-fas} 排除假解。
\end{remark}

\section{谱半径与「一次扫描的模态放大」：Gauss--Seidel 的两个数}
\label{sec:mg-gs}

\begin{lemma}[定常迭代的收敛判据（引自 1.5 章）]
\label{lem:mg-stationary}
对 $x^{(k+1)}=Tx^{(k)}+c$，迭代收敛（对任意初值）当且仅当 $\rho(T)<1$；
且误差按 $\rho(T)$ 渐近衰减。特别地本章的加权 Jacobi 收敛当且仅当
$\max_k|\mu_k(\omega)|<1$。
\end{lemma}

\begin{proof}
这是 1.5 章 的定常迭代判据与 Gelfand 公式的直接结论，这里只把结论搬过来：
$T^k\to0$ 当且仅当 $\rho(T)<1$（Jordan 块 + $\rho^k$ 因子），
而 $e^{(k)}=T^ke^{(0)}$。本章只需把它用在 $T_\omega$
（特征值 $\mu_k$，定理~\ref{thm:mg-spectrum}）上。
\end{proof}

\begin{theorem}[Gauss--Seidel 的无穷网格符号]
\label{thm:mg-gs-symbol}
在无穷网格（无边界）上，$T_{\mathrm{GS}}=(D-L)^{-1}U=(2I-S_-)^{-1}S_+$
（$S_\pm$ 是移位算子，$h^{-2}$ 因子完全约掉，符号与 $h$ 无关）。
把形式模态 $u_j=z^j$ 代入 $T_{\mathrm{GS}}u=\lambda u$ 得
\begin{equation}
  \lambda(z)=\frac{z^2}{2z-1},
  \qquad
  z=e^{\mathrm i\theta}\ \Rightarrow\
  |\lambda|=\frac{1}{|2e^{\mathrm i\theta}-1|}=\frac{1}{\sqrt{5-4\cos\theta}} .
  \label{eq:mg-gs-symbol}
\end{equation}
它在高频半带 $\theta\in[\pi/2,\pi]$ 上的最大值在 $\theta=\pi/2$ 取到，等于
\begin{equation}
  \max_{\theta\ge\pi/2}|\lambda|=\frac{1}{\sqrt5}=0.4472\ldots
  \label{eq:mg-gs-max}
\end{equation}
（端点 $\theta=\pi$ 给出 $1/3$）。
\end{theorem}

\begin{proof}
$A=D-L-U$，其中 $D=2h^{-2}I$、$L=h^{-2}S_-$、$U=h^{-2}S_+$。
于是 $(D-L)^{-1}U=(2I-S_-)^{-1}S_+$，$h^{-2}$ 被约掉。
设 $u_j=z^j$：$(S_+u)_j=z^{j+1}$、$(S_-u)_j=z^{j-1}$，代入 $S_+u=\lambda(2I-S_-)u$ 得
\begin{equation}
  z^{j+1}=\lambda\bigl(2z^j-z^{j-1}\bigr)
  \ \Longrightarrow\
  z=\lambda\bigl(2-z^{-1}\bigr)
  \ \Longrightarrow\
  \lambda=\frac{z^2}{2z-1}.
\end{equation}
取 $z=e^{\mathrm i\theta}$：$|2e^{\mathrm i\theta}-1|^2=5-4\cos\theta$，
故 $|\lambda|=(5-4\cos\theta)^{-1/2}$。
在 $[\pi/2,\pi]$ 上 $\cos\theta$ 从 $0$ 单调降到 $-1$，所以 $5-4\cos\theta$ 从 $5$ 单调升到 $9$，
最小值 $5$ 在 $\theta=\pi/2$ 取到，故最大值是 $1/\sqrt5=0.4472\ldots$，端点 $\theta=\pi$ 是 $1/3$。
\end{proof}

\begin{remark}[两个不同的数，别混用]
对同一维模型问题，Gauss--Seidel 有两个都很重要的数字：
\begin{itemize}
  \item \textbf{谱半径}（决定作为求解器的渐近速率）：
        由 1.5 章的 Young 关系，特征值是 $\cos^2\theta_k$，最大的是
        $\cos^2(\pi/(n+1))=1-O(h^2)$——低频几乎不动，与 \eqref{eq:mg-lowmode} 同一量级。
  \item \textbf{一次扫描对一个振荡分量的放大}：定理~\ref{thm:mg-gs-symbol} 的符号
        $1/\sqrt{5-4\cos\theta}$，在 $\theta=\pi/2$ 附近约为 $0.4472$——
        高频一次就掉一半多。
\end{itemize}
两者差这么远，是因为 $T_{\mathrm{GS}}$ **不是正规矩阵**：
它的特征向量（低频形状的伸缩模态）与它实际放大/衰减的模态（振荡的移位模态）不是同一组，
谱半径只描述最终速率，不描述一次扫描对某个振荡分量的即时效果。
多重网格的平滑性分析用的正是后者（定理~\ref{thm:mg-smooth} 对 Jacobi 是精确的，
对 GS 则用符号 \eqref{eq:mg-gs-max}），这也是为什么 $\omega=2/3$ 的加权 Jacobi
在多重网格里比 GS 更常被拿来做分析：它的迭代矩阵在正弦基下**对角**，
逐模态放大率是精确的，而 GS 的同一句话只能作为内点的渐近描述。
\end{remark}

\section{本章代码与结论的对应}
\label{sec:mg-code}

\begin{itemize}
  \item \texttt{mode\_factors\_theory} / \texttt{mode\_factors\_measured}
        $\leftrightarrow$ 定理~\ref{thm:mg-spectrum}：前者给闭式
        $|\mu_k(\omega)|$，后者给算子作用下的实测比值；
        对 Jacobi 两者逐模态相符，对 GS 会出现模态混合（\S\ref{sec:mg-gs}）。
  \item \texttt{smoothing\_factor} $\leftrightarrow$ 定理~\ref{thm:mg-smooth}：
        $\omega=2/3$ 时高频半带上的最坏放大率 $1/3$。
  \item \texttt{galerkin\_defect} / \texttt{coarse\_scaling\_ratio}
        $\leftrightarrow$ 定理~\ref{thm:mg-galerkin}、推论~\ref{cor:mg-scaling}：
        缺陷恒为 $0$，比例恒为 $2$（$P^{\mathsf T}A_hP=2A_H$）。
        \texttt{restriction\_matrix} 的 \texttt{factor=0.5} 不是可调参数，
        它被恒等式定死；换成别的值就落到推论~\ref{cor:mg-scaling} 的
        $RA_hP=2cA_H$，粗层校正要么欠、要么过（$c=1$ 时过一倍）。
  \item \texttt{two\_grid\_matrix} / \texttt{two\_grid\_factor}
        $\leftrightarrow$ 定理~\ref{thm:mg-tg-spectrum}：谱是 $\{0,1/9\}$，
        幂迭代量到的就是 $\rho=1/9$。
  \item \texttt{mg\_hierarchy} / \texttt{mg\_apply} / \texttt{cycle\_factor} /
        \texttt{factor\_scan} $\leftrightarrow$ 定义~\ref{def:mg-cycle} 与
        定理~\ref{thm:mg-cost}：递归算子与它的代价计数器（点扫描数）；
        \texttt{factor\_scan} 把因子沿 $h$ 序列扫一遍，用来看「网格无关」是否真的成立。
  \item \texttt{MGPrecond} $\leftrightarrow$ 定理~\ref{thm:mg-cycle-spd}、
        \ref{thm:mg-mgcg}：零初值一次循环就是预条件子 $B$ 的作用。
  \item \texttt{nonlinear\_smooth} / \texttt{fas\_cycle} / \texttt{measure\_fas\_factor} /
        \texttt{discrete\_newton} $\leftrightarrow$ 定理~\ref{thm:mg-cubic}、
        \ref{thm:mg-fas}、\ref{thm:mg-false}：
        逐点三次方程用二分（区间由 \eqref{eq:mg-cubic-root} 的括住性质给出），
        FAS 的粗层方程是 \eqref{eq:mg-fas-coarse}，
        牛顿法（自研 LU）提供独立参考解，用来区分「残差小」与「解对」。
        \texttt{measure\_fas\_factor} 的 \texttt{nonlinear=False} 分支就是
        定理~\ref{thm:mg-false} 的反例入口。
  \item 合规说明（P2/P3）：本章不调用任何现成求解器。
        粗层求解用的是自研追赶法（\texttt{numcore.linalg.lu} 的三对角 LU 族），
        因子用自研幂迭代测，随机初值只经 \texttt{numcore.utils.rng}（P4），
        牛顿参考解也用自研 LU。
\end{itemize}

\section{它会怎样骗你（本章版）}

\begin{enumerate}
  \item \textbf{把「多重网格一步到位」当成默认}：V-cycle 的实测收敛因子是
        $0.1837$（$n=15$）、$0.1894$（$n=31$）、$0.1887$（$n=63$）、$0.1902$（$n=127$），
        降到 $10^{-8}$ 要\textbf{真的数}到 $11$ 个周期。它不是 $0$：固定递归的循环只是一个
        收缩算子，粗层永远只解了个近似。
  \item \textbf{把「网格无关」说成「因子是常数」}：同一条 $h$ 序列上因子是
        $0.1609/0.1837/0.1894/0.1887/0.1902$，均值 $0.1826$、标准差 $1.1\times10^{-2}$、
        漂移 $+2.9\times10^{-2}$——它是\textbf{缓慢变大}的。$n=7$ 那个 $0.1609$ 还是小网格异常：
        最粗层只有 $1$ 个内点却做了 $24$ 次扫描，等于把粗层过度求解了。
  \item \textbf{拿平滑因子给循环因子下结论}：Gauss--Seidel 的高频平滑因子实测
        $0.4453$（$n=127$，趋近 $1/\sqrt5=0.4472$），比最优 Jacobi 的 $1/3$ \emph{更差}；
        可同一个循环里 GS 的 V-cycle 因子是 $0.104$（$n=31$）与 $0.134$（$n=127$），
        比 Jacobi 的 $0.190$ \emph{更好}。原因是低频：$\rho(T_{\rm GS})=\cos^2\theta_1$，
        比 Jacobi 的 $\cos\theta_1$ 快一倍。
  \item \textbf{把理论曲线当实际衰减}：$n=127$、$\omega=1.2$ 的 SOR，迭代矩阵特征值口径给
        $0.999097$（几乎不动），而单模态一步范数比只有 $0.541905$。SOR 的迭代矩阵不是正规矩阵，
        特征值模长只决定渐近速率；两条线相差 $0.46$，不是画错。
  \item \textbf{只看整体误差范数}：$n=63$、初值 $c_k=1/k$、$10$ 步之后整段范数只掉到
        $0.936$，看上去「平滑几乎没用」；同一时刻高频模态从 $1.587\times10^{-2}$ 掉到
        $1.067\times10^{-9}$，低/高存活比 $6.0\times10^{4}$。范数那条平线完全由
        $\lambda_1=0.999197$ 决定。
  \item \textbf{限制算子随便定个缩放}：粗层算子固定为重离散、$R=cP^{\mathsf T}$ 时，
        $c=0.25$ 的因子是 $0.8919$（像「只是慢一点」），$c=0.75$ 就跳到 $2.0433$，
        $c=1.0$ 是 $6.0376$，$c=2.0$ 是 $53.4368$——循环\textbf{发散}。
        Galerkin 缺口 $\max|RA_hP-A_H|$ 恰好等于 $|2c-1|$，是最省事的体检量。
  \item \textbf{以为 injection 与满权只差一个常数}：injection 只取 $r_{2j+1}$，
        破坏的是伴随关系本身。两网格因子从 $0.111111$ 退到 $0.332263$（$n=63$），
        而且把最优常数 $c$ 调到最好，$\|R_{\rm inj}-cP^{\mathsf T}\|$ 仍有 $0.5$ 量级。
  \item \textbf{只看因子不看代价}：$n=63$ 上 W-cycle 的因子是 $0.1145$，确实比 V 的 $0.1887$ 小
        （接近两网格的 $1/9$）；但每个周期的实测代价是 $21.365$ 个「最细层等价点扫描」，
        V 只有 $4.159$。降到 $10^{-8}$ 的总代价是 $9\times21.365=192.3$ 对 $11\times4.159=45.7$
        ——W 贵 $4.2$ 倍。「因子更小」被「一步更贵」吃掉。
  \item \textbf{把加速比说成时间比}：同一右端、$\mathrm{tol}=10^{-8}$，CG 的迭代数恰好是
        $n$（$15/31/63/127/255$），MGCG 恒为 $7$，迭代数加速比
        $2.14/4.43/9.00/18.14/36.43$。但 MGCG 每次迭代要多做一个 V-cycle（约 $4$ 个
        「最细层等价代价」），$n=15$ 时总代价只比单独跑 V-cycle 好 $1.6$ 倍。
        「加速 $36$ 倍」必须紧跟「这是 $n=255$ 的迭代数比」。
  \item \textbf{顺手写下的两个「显然」恒等式}（本节作者真踩过）：第一版断言里写了
        $RA_h^{-1}=2A_H^{-1}R$，一般右端上相对偏差 $1.0098$；退到子空间
        $A_h\,\mathrm{range}(P)$ 上也差 $0.6126$。更隐蔽的是 $RP=I$：
        一维里 $RP=0.75I+0.125\,\mathrm{adj}$，$\max|RP-I|=0.25$
        （因为 $P^{\mathsf T}P\ne2I$）。真命题是 $A_H^{-1}RA_hP=I$ 与
        $Q=PA_H^{-1}RA_h$ 幂等、$A_h$-对称。
  \item \textbf{因子估计器读错一位}（同一个坑的另一半）：把「每步归一化后的范数」再相邻相除，
        $n=15$ 上给出 $1.0038$，看起来像循环不收敛；显式循环矩阵的谱半径是 $0.18955$，
        未归一化迭代 $1.24\times10^{-1}\to7.93\times10^{-10}$。归一化后的范数\textbf{就是}
        那一步的放大率，尾部几何平均才是 $\rho$。
  \item \textbf{把残差的舍入地板当收敛停滞}：FAS 在 $n=63$、$\nu=2$ 时残差很快压到
        $10^{-11}$ 附近，此时相邻比值全是 $1\pm$ 噪声，拿最后四个比值求平均得到 $0.8566$
        （真值 $0.0774$）。因子只能在\textbf{触底之前}的窗口里估。
  \item \textbf{把非线性问题当线性解}：$n=63$ 的 $-u''+u^3=f$ 上，丢掉粗层的非线性项以后
        \emph{线性}残差压到 $4.075\times10^{-10}$（按它自己的尺子「收敛了」），
        真残差却停在 $5.5960$，解偏了 $7.886\times10^{-2}$（真解幅度 $1$，即 $8\%$ 误差），
        两个残差的比值 $1.4\times10^{10}$。FAS 的正常路径同一算例的真残差是
        $3.07\times10^{-10}$，与离散牛顿参考解差 $1.3\times10^{-13}$。
  \item \textbf{以为粗网格「看不见」高频就是「丢掉了」高频}：$\theta/\pi=0.9531$ 的纯高频残差，
        限制后只剩原幅度的 $3.827\times10^{-3}$，别名到粗网格频率 $-0.0938\pi$（低频），
        而粗层可见部分 $PRr$ 的高频能量比是 $0.0000$。粗网格上的一切都是低频，
        包括那些本来是高频的东西。
  \item \textbf{以为多平滑几步高频就「变得可见」了}：同一个纯高频残差，$\nu=0,\dots,6$ 的
        相对损失恒为 $0.999985$（形状不变），只有幅度每步乘 $0.3261$；真正会变的是
        \emph{混合}残差：高频能量占一半时相对损失从 $0.6398$ 掉到 $0.0018$。
        平滑改变的是高频的\textbf{大小}，不是它的\textbf{可见性}。
  \item \textbf{把 $O(h^2)$ 当斜率读数}：$1-\rho(T_J)=\pi^2h^2/2+O(h^4)$ 的渐近斜率是 $2$，
        实测（$n=7,\dots,127$）是 $-1.889$——最小那个网格还没进渐近区。同理，
        单层平滑子降到 $10^{-8}$ 需要 $6.4\times10^{4}$ 次扫描（$n=127$，实测因子 $0.999710$），
        帧里的这一列是\textbf{换算值}，MG 那一列（$11$ 个周期）才是真的数出来的。
  \item \textbf{把一维的「精确」搬去二维}：$RA_hP=A_H$ 在 1D 逐位成立，所以本章可以放心用
        重离散粗层算子；二维双线性插值配满权限制给出的是 $9$ 点 Galerkin 模板，
        \emph{不是} $5$ 点重离散，两者差一个 $O(h^2)$ 的交叉项。本章只做 1D，
        这一条不许当前提用。
\end{enumerate}

% 本文件由 tools/build_book.py 从 dist/reports/assertions.json 生成，请勿手改。
\begin{longtable}{@{}p{0.06\textwidth}p{0.32\textwidth}p{0.22\textwidth}p{0.16\textwidth}@{}}
\caption{章节 4.5 数值断言（生成时间 2026-10-18T17:31:25-0400）}\\
\toprule
类别 & 断言 & 理论值 & 实测值 \\
\midrule
\endfirsthead
\toprule
类别 & 断言 & 理论值 & 实测值 \\
\midrule
\endhead
\bottomrule
\endlastfoot

B & 帧角点：有限性 + requiredAux + 掩码自洽 & 数组与 meta 全有限、\allowbreak{}能编 .nlf；requiredAux 命中；地板/\allowbreak{}占位/\allowbreak{}哨兵掩码与数据一致 & \{'first': [],\allowbreak{} 'n\_frames': 104,\allowbreak{} 'n\_problems': 0\} \checkmark \\
B & meta 与数组的自洽性(四页) & meta 的每个结论都能从同帧数组复算(\allowbreak{}平滑因子/\allowbreak{}地板计数/\allowbreak{}因子统计/\allowbreak{}损失与别名) & \{'n\_problems': 0,\allowbreak{} 'problems': []\} \checkmark \\
B & 纯函数与时间钳制 & 同参数两次构建逐位相同；scrub 在 [0,\allowbreak{}tMax] 外钳制；mode=\allowbreak{}none 不声明 tMax/\allowbreak{}dt/\allowbreak{}fps & \{'n\_problems': 0,\allowbreak{} 'problems': []\} \checkmark \\
C & 跨页阈值：高频消失、低频不动、MG 与平滑子的量级差 & 10 步后高频 <=\allowbreak{}1e-4 而低频 >=\allowbreak{}0.99；MG 总代价比平滑子 >=\allowbreak{}100 倍；纯高频残差的不可见比例 >=\allowbreak{}0.9 且混合残差损失至少掉 10 倍 & \{'viz1\_high\_survival': 1.686706566226433e-05,\allowbreak{} 'viz1\_low\_survival': 0.9979938376613837,\allowbreak{} 'viz3\_cost\_ratio': 800.3068181818182,\allowbreak{} 'viz3\_mgcg\_iters': 7.0,\allowbreak{} 'viz3\_two\_grid': 0.11111111111111113,\allowbreak{} 'viz4\_mixed\_drop': 390.82270979700394,\allowbreak{} 'viz4\_pure\_loss': 0.9999851987624295,\allowbreak{} 'viz4\_pure\_norm': 0.0012029456640646609\} \checkmark \\
B & manifest 自洽性 & 参数边界/\allowbreak{}默认值/\allowbreak{}临界值/\allowbreak{}枚举自洽，轴 >=\allowbreak{} 2 且 autoscaled 显式，pitfalls 与 theory 非空 & \{'n\_problems': 0,\allowbreak{} 'problems': []\} \checkmark \\
A & 转移恒等式最大偏差 & max|R-0.5P\textasciicircum{}T|=\allowbreak{}0，max|R A\_h P-A\_H|=\allowbreak{}0，max|P\textasciicircum{}T A\_h P-2A\_H|=\allowbreak{}0 & \{'PtAP\_minus\_2AH': 0.0,\allowbreak{} 'R\_minus\_half\_PT': 0.0,\allowbreak{} 'galerkin\_defect': 0.0,\allowbreak{} 'poisson\_vs\_1\_5': 0.0\} \checkmark \\
A & Jacobi 逐模态因子：实测 vs 理论 vs 1.5 章口径 & |1-omega(\allowbreak{}1-cos theta\_k)|；omega=\allowbreak{}2/\allowbreak{}3 时 |1-(\allowbreak{}4/\allowbreak{}3)sin\textasciicircum{}2(\allowbreak{}theta\_k/\allowbreak{}2)| & \{'measured\_vs\_theory': 4.447949739383039e-15,\allowbreak{} 'theory\_vs\_closed\_form': 2.220446049250313e-16\} \checkmark \\
A & smoother\_\allowbreak{}matrix vs 1.5 章迭代矩阵 & 与 jacobi/\allowbreak{}gs/\allowbreak{}sor\_iteration\_matrix 逐位相同(\allowbreak{}omega 由 I-omega(\allowbreak{}I-T\_J) 缩放) & 0 \checkmark \\
A & 两网格因子 vs 1/9 & rho\_TG =\allowbreak{} mu\_smooth\textasciicircum{}2 =\allowbreak{} (\allowbreak{}1/\allowbreak{}3)\textasciicircum{}2 =\allowbreak{} 0.111111... & \{'factors': [0.11111111111111104,\allowbreak{} 0.11111111111111067,\allowbreak{} 0.11111111111111116,\allowbreak{} 0.11111111111111077],\allowbreak{} 'max\_abs\_dev': 4.3021142204224816e-16,\allowbreak{} 'spread': 4.85722573273506e-16\} \checkmark \\
A & FAS vs 离散牛顿参考解 & FAS 不动点 =\allowbreak{} 离散方程的解：max|u\_FAS - u\_Newt| 到机器精度；残差 <=\allowbreak{} 1e-8 & \{'discrete\_vs\_exact': [0.0025975536252749176,\allowbreak{} 0.0006490957336637848,\allowbreak{} 0.00016225547489456815],\allowbreak{} 'exact\_error': [7.593925488436071e-14,\allowbreak{} 1.156852391659413e-13,\allowbreak{} 1.3233858453531866e-13],\allowbreak{} 'final\_residual': [1.3547766369373269e-11,\allowbreak{} 7.5482179359868e-11,\allowbreak{} 3.071030313191496e-10]\} \checkmark \\
B & 粗层校正对 P 值域的零化(相对) & R A\_h P =\allowbreak{} A\_H $\Rightarrow$ (\allowbreak{}I-P A\_H\textasciicircum{}\{-1\} R A\_h) P y =\allowbreak{} 0 & $3.082\times 10^{-14}$ \checkmark \\
B & Galerkin 投影 vs 两个假恒等式 & Q\textasciicircum{}2=\allowbreak{}Q、\allowbreak{}A-对称(\allowbreak{}偏差 0)；R P =\allowbreak{} 0.75I+0.125 adj(\allowbreak{}偏差 0)；max|RP-I| =\allowbreak{} 0.25 !=\allowbreak{} 0 & \{'projection\_defect': 4.4408920985006214e-15,\allowbreak{} 'rp\_closed\_form': 0.0,\allowbreak{} 'rp\_vs\_identity': 0.25\} \checkmark \\
B & 含逆交换律的反例留档 & 一般右端 rel$\approx$1.0；子空间上 abs$\approx$0.61；真命题 A\_H\textasciicircum{}\{-1\} R A\_h P =\allowbreak{} I 偏差 0 & \{'general\_rhs\_rel': 1.0098227246529399,\allowbreak{} 'subspace\_abs': 0.6126289540840444,\allowbreak{} 'true\_form\_abs': 1.7763568394002505e-15\} \checkmark \\
B & 循环算子线性性与矩阵一致性(相对) & E 线性：E(\allowbreak{}x1+x2,\allowbreak{}b1+b2)=\allowbreak{}E(\allowbreak{}x1,\allowbreak{}b1)+E(\allowbreak{}x2,\allowbreak{}b2)，且 E 的列就是 mg\_apply(\allowbreak{}e\_j) & $7.406\times 10^{-15}$ \checkmark \\
B & 计数器 vs 闭式代价(相对) & V: sum\_l 2*n\_l/\allowbreak{}n\_0 + 24*n\_L/\allowbreak{}n\_0；W: 重数 2\textasciicircum{}l & \{'rows': [[15,\allowbreak{} 'v',\allowbreak{} 4.933333333333334,\allowbreak{} 4.933333333333334],\allowbreak{} [15,\allowbreak{} 'w',\allowbreak{} 18.26666666666667,\allowbreak{} 18.266666666666666],\allowbreak{} [31,\allowbreak{} 'v',\allowbreak{} 4.387096774193548,\allowbreak{} 4.387096774193548],\allowbreak{} [31,\allowbreak{} 'w',\allowbreak{} 19.677419354838705,\allowbreak{} 19.67741935483871]],\allowbreak{} 'worst\_rel': 4.988580128479158e-16\} \checkmark \\
B & 估计器：对法 vs 错法 vs 显式矩阵谱半径 & 归一化后每步范数即放大率(\allowbreak{}几何平均 -> rho(\allowbreak{}E))；相邻相除会得到 $\approx$1 的假值 & \{'cycle\_factor\_api': 0.1837371257461596,\allowbreak{} 'cycle\_factor\_correct': 0.1837371257461596,\allowbreak{} 'naive\_ratio\_of\_ratios': 1.0037651335386966,\allowbreak{} 'rho\_of\_matrix': 0.18955483498701953\} \checkmark \\
B & injection 的伴随缺口与两网格因子代价 & R =\allowbreak{} 0.5 P\textasciicircum{}T(\allowbreak{}满权)才使 R A\_h P =\allowbreak{} A\_H；injection 下两网格因子显著变大(\allowbreak{}实测约 3 倍) & \{'fw': [0.11111111111111067,\allowbreak{} 0.11111111111111116,\allowbreak{} 0.11111111111111077],\allowbreak{} 'inj': [0.32906339564483605,\allowbreak{} 0.3322632725939292,\allowbreak{} 0.33306559645892087],\allowbreak{} 'min\_ratio': 2.961570560803536,\allowbreak{} 'worst\_asymmetry': 0.33333333333333337\} \checkmark \\
C & GS 高频放大率与边界误差收缩 & max\_high |g(\allowbreak{}theta)| -> 1/\allowbreak{}sqrt(\allowbreak{}5)=\allowbreak{}0.447214；边界误差 \textasciitilde{} O(\allowbreak{}h)(\allowbreak{}比值 -> 2) & \{'gap': [0.024552769515947204,\allowbreak{} 0.01224299769295889,\allowbreak{} 0.006093617540537721,\allowbreak{} 0.0030413549001271356],\allowbreak{} 'gap\_ratios': [2.0054540670270504,\allowbreak{} 2.0091509865055506,\allowbreak{} 2.0035864740030815],\allowbreak{} 'max\_abs\_dev': 0.0018672883037116894,\allowbreak{} 'max\_high': [0.43205008623077534,\allowbreak{} 0.4396968652810592,\allowbreak{} 0.443471156521669,\allowbreak{} 0.44534630719624624],\allowbreak{} 'young\_eigen': [0.9619397662556434,\allowbreak{} 0.9903926402016151,\allowbreak{} 0.9975923633360985,\allowbreak{} 0.9993977281025862]\} \checkmark \\
C & 缩放常数 c 的因子扫描 & |c-0.5|=\allowbreak{}0 时循环收敛；c>0.5 过校正可发散(\allowbreak{}粗层校正被放大 2c 倍) & \{'c': [0.25,\allowbreak{} 0.5,\allowbreak{} 0.75,\allowbreak{} 1.0,\allowbreak{} 2.0],\allowbreak{} 'cycle\_factor': [0.9345419144624127,\allowbreak{} 0.188663987137861,\allowbreak{} 4.343326780434276,\allowbreak{} 19.196681369289028,\allowbreak{} 549.1505837323252],\allowbreak{} 'galerkin\_defect': [0.5,\allowbreak{} 0.0,\allowbreak{} 0.5,\allowbreak{} 1.0,\allowbreak{} 3.0],\allowbreak{} 'worst': 549.1505837323252\} \checkmark \\
C & V vs W：因子、每周期代价、总代价 & W 因子接近两网格(\allowbreak{}更小)，但每周期代价约 2\textasciicircum{}L 级；总代价由 因子 与 代价 共同决定 & \{'cost\_ratio': 4.203331020124932,\allowbreak{} 'cost\_v': 45.74603174603173,\allowbreak{} 'cost\_w': 192.2857142857151,\allowbreak{} 'v\_cycles': 11,\allowbreak{} 'v\_factor': 0.188663987137861,\allowbreak{} 'v\_work': 4.158730158730157,\allowbreak{} 'w\_cycles': 9,\allowbreak{} 'w\_factor': 0.11448953348269528,\allowbreak{} 'w\_work': 21.365079365079456,\allowbreak{} 'work\_ratio': 5.137404580152696\} \checkmark \\
C & 因子扫描：MG 的平坦 vs 平滑子的 h\textasciicircum{}2 恶化 & MG 因子与 h 无关(\allowbreak{}std 小、\allowbreak{}漂移小)；1-rho\_smoother \textasciitilde{} h\textasciicircum{}2 $\Rightarrow$ 斜率 -2(\allowbreak{}渐近) & \{'drift': 0.02922987995958487,\allowbreak{} 'factor\_mg': [0.16093827456842116,\allowbreak{} 0.1837371257461596,\allowbreak{} 0.18935184490011966,\allowbreak{} 0.188663987137861,\allowbreak{} 0.19016815452800603],\allowbreak{} 'factor\_smoother': [0.9492530216741912,\allowbreak{} 0.9871901869343185,\allowbreak{} 0.9967898177764185,\allowbreak{} 0.9991909271994155,\allowbreak{} 0.9997104469939833],\allowbreak{} 'factor\_two\_grid': [0.11111111111111113,\allowbreak{} 0.11111111111111104,\allowbreak{} 0.11111111111111067,\allowbreak{} 0.11111111111111116,\allowbreak{} 0.11111111111111077],\allowbreak{} 'iters\_mg': [2.0,\allowbreak{} 2.0,\allowbreak{} 2.0,\allowbreak{} 2.0,\allowbreak{} 2.0],\allowbreak{} 'iters\_smoother': [56.0,\allowbreak{} 226.0,\allowbreak{} 907.0,\allowbreak{} 3600.0,\allowbreak{} 10061.0],\allowbreak{} 'mean': 0.18257187737611352,\allowbreak{} 'n': [7,\allowbreak{} 15,\allowbreak{} 31,\allowbreak{} 63,\allowbreak{} 127],\allowbreak{} 'smoother\_slope': -1.8891537362805895,\allowbreak{} 'std': 0.011046771127791241\} \checkmark \\
C & CG vs MGCG 迭代数(tol=1e-8) & CG \textasciitilde{} O(\allowbreak{}h\textasciicircum{}-1)(\allowbreak{}sqrt(\allowbreak{}kappa))；MGCG 迭代数 O(\allowbreak{}1)；加速比按迭代数 \textasciitilde{} 1/\allowbreak{}h & \{'cg\_iters': [15,\allowbreak{} 31,\allowbreak{} 63,\allowbreak{} 127,\allowbreak{} 255],\allowbreak{} 'mgcg\_iters': [7,\allowbreak{} 7,\allowbreak{} 7,\allowbreak{} 7,\allowbreak{} 7],\allowbreak{} 'mgcg\_std': 0.0,\allowbreak{} 'n': [15,\allowbreak{} 31,\allowbreak{} 63,\allowbreak{} 127,\allowbreak{} 255],\allowbreak{} 'speedup': [2.142857142857143,\allowbreak{} 4.428571428571429,\allowbreak{} 9.0,\allowbreak{} 18.142857142857142,\allowbreak{} 36.42857142857143]\} \checkmark \\
C & FAS 因子(跨 h)与线性化反例 & FAS 因子与 h 无关；线性化残差收敛到 1e-10 时真残差仍 O(\allowbreak{}1)，假解误差 \textasciitilde{}8e-2(\allowbreak{}真解幅度 1) & \{'false\_linear\_residual': 4.075381651270153e-10,\allowbreak{} 'false\_solution\_error': 0.07886428380016719,\allowbreak{} 'false\_true\_residual': 5.5964906255189675,\allowbreak{} 'fas\_factor': [0.15442316308010817,\allowbreak{} 0.1734676284798888,\allowbreak{} 0.18617952604063434,\allowbreak{} 0.18917931305791433],\allowbreak{} 'gap': 13732433191.31777,\allowbreak{} 'mean': 0.17581240766463643,\allowbreak{} 'n': [15,\allowbreak{} 31,\allowbreak{} 63,\allowbreak{} 127],\allowbreak{} 'std': 0.013685293900176134\} \checkmark \\
B & 窗口因子 vs 触底后的朴素比值 & 触底段比值 \textasciitilde{}1；窗口估计器跳过地板与初始瞬态后应给出真衰减率 0.19 & \{'naive\_tail\_ratio': 1.0,\allowbreak{} 'window': [7,\allowbreak{} 11],\allowbreak{} 'window\_factor': 0.19\} \checkmark \\
\end{longtable}

